% ID: SEL-D0007-BA-E02
% ESTADO: cotejada. Segunda pasada visual completa y borrador compilado revisados.
% Fuente: Selectividad/matematicas_aplicadas/ORD_TITULAR_MATEMÁTICAS APLICADAS CCSS-JUNIO-IMPRESO.pdf
% SHA-256: c42a605501a277fe7194b964ad92ad66047be9a62dc254fa2219d756d3d71317
% Página: 1; bloque A; ejercicio original 2.

Se considera la matriz
$A=\begin{pmatrix}1&-1&m\\0&2&-3\\m&1&1\end{pmatrix}$,
con $m$ un parámetro real.
\begin{enumerate}[label=\alph*)]
  \item \textbf{(0.7 puntos)} \textquestiondown{}Para qué valores del
  parámetro $m$ existe la matriz inversa de $A$?
  \item \textbf{(1.8 puntos)} Para $m=2$, resuelva la ecuación matricial
  $X\cdot A-A^2=I_3$.
\end{enumerate}
